\begin{abstract}
Scientific papers contain valuable metadata about experimental facilities,
samples, measurement datasets, and their relationships, but this information
is distributed across the publication rather than reported in a single
structured form. We introduce a publication-derived benchmark for structured
metadata extraction from 40 photon and neutron (PaN) science papers, with
reconciled gold records under a two-layer schema covering generic
publication and PaN experimental metadata. To account for repeated entities
and cross-entity references expressed through output-local identifiers, 
we develop an entity-aware evaluation protocol that aligns predicted and 
gold entities, scores attributes on aligned pairs, and remaps local identifiers 
before scoring references.

We evaluate seven large language models (LLMs) on the 30-paper evaluation set in
zero-shot and instruction-ablation settings, with additional within-family scaling and
few-shot analyses. Generic metadata are recovered more consistently than
facilities, samples, and datasets, whose annotation disagreements frequently
concern entity identity and granularity. Detailed field instructions yield higher 
sample and dataset recovery across all models, while the effects of model size 
and few-shot demonstrations are less uniform. Higher entity recovery does not 
necessarily correspond to higher conditional attribute or reference quality.
\end{abstract}
